\documentclass[10pt]{article}
\usepackage[T1]{fontenc}
\usepackage[utf8]{inputenc}
\usepackage[french]{babel}
\usepackage{amsmath,amssymb,mathrsfs,mathpazo}
\usepackage{geometry}
\geometry{a4paper,margin=2cm}
\usepackage{graphicx,array,tabularx,tabularray,colortbl,multirow,multicol,booktabs,enumitem,pifont,ulem,fancybox,diagbox,cancel}
\usepackage{pstricks,pstricks-add,pst-plot,pst-tree,pst-node,pst-text,pst-func,pst-math,pst-bspline,pst-3dplot,pst-eucl,pst-solides3d}
\let\pscalculate\undefined
\usepackage{pst-bezier}
\usepackage{tikz,pgfplots,tkz-tab}
\usetikzlibrary{arrows,arrows.meta,patterns,calc,positioning,trees,decorations.markings,shapes.geometric,shapes.misc,shapes.arrows,intersections,decorations.pathmorphing,decorations.pathreplacing,backgrounds,shadows,babel}
\pgfplotsset{compat=1.18}
\usepgfplotslibrary{fillbetween}
\usepackage[np]{numprint}
\usepackage{listings,makecell,pas-tableur,twoopt}
\let\barre\undefined
\usepackage[pstricks]{tablvar}
\definecolor{darkgreen}{RGB}{0,100,0}
\providecommand{\red}{\color{red}}
\providecommand{\blue}{\color{blue}}
\providecommand{\green}{\color{green}}
\providecommand{\black}{\color{black}}
\providecommand{\white}{\color{white}}
\providecommand{\gray}{\color{gray}}
\providecommand{\cyan}{\color{cyan}}
\providecommand{\magenta}{\color{magenta}}
\providecommand{\yellow}{\color{yellow}}
\providecommand{\violet}{\color{violet}}
\providecommand{\orange}{\color{orange}}
\providecommand{\euro}{\texteuro}
\providecommand{\EUR}{\texteuro}
\providecommand{\ieme}{\textsuperscript{e}}
\providecommand{\degres}{\ensuremath{^\circ}}
\providecommand{\degree}{\ensuremath{^\circ}}
\providecommand{\vv}[1]{\overrightarrow{#1}}
\providecommand{\Coord}[1]{#1}
\providecommand{\decosix}{}
\providecommand{\starredbullet}{\ensuremath{\bullet}}
\pagestyle{empty}
\setlength{\parindent}{0pt}
\usepackage{adjustbox}\DeclareUnicodeCharacter{2236}{:}\DeclareUnicodeCharacter{2190}{\ensuremath{\leftarrow}}
\providecommand{\R}{\mathbb{R}}
\providecommand{\N}{\mathbb{N}}
\providecommand{\Z}{\mathbb{Z}}
\providecommand{\Q}{\mathbb{Q}}
\providecommand{\D}{\mathbb{D}}
\providecommand{\C}{\mathbb{C}}
\providecommand{\e}{{\mathrm{e}}}
\providecommand{\vect}[1]{\overrightarrow{#1}}
\providecommand{\vectt}[1]{\overrightarrow{\mathrm{#1}}}
\providecommand{\vvt}[1]{\overrightarrow{\mathrm{#1}}}
\providecommand{\np}[2][]{{#2~#1}}
\providecommand{\vv}[1]{{\overrightarrow{#1}}}
\renewcommand{\d}{{\mathrm{d}}}
\providecommand{\degres}{{^\circ}}
\providecommand{\degre}{{^\circ}}
\providecommand{\blue}{\color{blue}}
\providecommand{\red}{\color{red}}
\providecommand{\enskip}{\quad}
\providecommand{\llbracket}{\lbrack\!\lbrack}
\providecommand{\rrbracket}{\rbrack\!\rbrack}
\providecommand{\euro}{{\text{€}}}
\providecommand{\no}{{\text{n°}}}
\providecommand{\interval}[2]{{[#1;#2]}}
\providecommand{\intervFO}[2]{{[#1;#2[}}
\providecommand{\dotfill}{\ldots}

\providecommand{\ds}{\displaystyle}
\providecommand{\textding}[1]{#1}
\providecommand{\Oij}{\(\mathrm{O};\vec{\imath},\vec{\jmath}\)}
\providecommand{\Oijk}{\(\mathrm{O};\vec{\imath},\vec{\jmath},\vec{k}\)}
\providecommand{\Ouv}{\(\mathrm{O};\vec{u},\vec{v}\)}
\providecommand{\intervFF}[2]{\ensuremath{[#1;#2]}}
\renewcommand{\euro}{\ifmmode\text{\texteuro}\else\texteuro\fi}
\DeclareUnicodeCharacter{00B0}{\ifmmode^{\circ}\else\textdegree\fi}
\DeclareUnicodeCharacter{20AC}{\ifmmode\text{\texteuro}\else\texteuro\fi}
% Énoncés : Éducation nationale. Transcription et corrigés : APMEP, Denis Vergès et contributeurs. Découpage et classement : CoopMaths, Sébastien Lozano, Éric Elter et contributeurs. Adaptation de la présentation : Maths Explorer.
% CC BY-SA : https://coopmaths.fr/www/about/
% https://coopmaths.fr/alea/static/bac/2024/tex/bac_2024_09_sujet1_metropole_1.tex
\begin{document}
\noindent\textit{APMEP / CoopMaths. CC BY-SA. Présentation adaptée par Maths Explorer.}\par\medskip
\medskip

On considère un cube ABCDEFGH de côté 1.

\begin{center}
\psset{unit=0.8cm}
\begin{pspicture}(9,9.4)
\psdots(1.4,1)(5.6,0.5)(7.9,1.4)(3.7,1.9)(1.4,5.5)(5.6,5)(7.9,5.9)(3.7,6.4)(4.65,1.2)(1.4,8.2)(7.1,4.72)
\pspolygon(1.4,1)(5.6,0.5)(5.6,5)(1.4,5.5)
\psline(5.6,0.5)(7.9,1.4)(7.9,5.9)(5.6,5)
\psline(7.9,5.9)(3.7,6.4)(1.4,5.5)
\psline(1.4,8.2)(5.6,0.5)(7.9,5.9)
\psline(1.4,8.2)(5.5,5.7)
\psline[linestyle=dashed](5.5,5.7)(7.9,4.2)
\psline(1.4,8.2)(2.24,5.98)
\psline(1.4,1)(1.4,9)
\psline(0.4,8.8)(1.4,8.2)
\psline[linestyle=dashed](2.24,5.98)(3.7,1.9)
\psline(7.9,4.2)(8.7,3.7)
\psline[linestyle=dashed](1.4,1)(3.7,1.9)(7.9,1.4)
\psline[linestyle=dashed](3.7,6.4)(3.7,1.9)(5.6,0.5)
\psline[linestyle=dashed](4.65,1.2)(7.9,5.9)(3.7,1.9)
\uput[l](1.4,1){A} \uput[dr](5.6,0.5){B} \uput[r](7.9,1.4){C} \uput[170](3.7,1.9){D}
\uput[l](1.4,5.5){E} \uput[ul](5.6,5){F} \uput[ur](7.9,5.9){G} \uput[170](3.7,6.4){H}
\uput[dl](4.65,1.2){I} \uput[ur](1.4,8.2){K} \uput[d](7.1,4.72){L} \uput[dl](0.4,8.8){$\Delta$}
\end{pspicture}
\end{center}

Le point I est le milieu du segment [BD]. On définit le point L tel que $\vect{\text{IL}}= \dfrac34\vect{\text{IG}}$.

On se place dans le repère orthonormé $\left(\text{A}~;~\vect{\text{AB}},\:\vect{\text{AD}},\:\vect{\text{AE}}\right)$.

\begin{enumerate}
\item
	\begin{enumerate}
		\item Préciser les coordonnées des points D, B, I et G. \\
		Aucune justification n'est attendue.
		\item Montrer que le point L a pour coordonnées $\left(\dfrac78~;~\dfrac78~;~\dfrac34\right)$.
	\end{enumerate}
\item Vérifier qu'une équation cartésienne du plan (BDG) est $x + y - z - 1 = 0$.
\item On considère la droite $\Delta$ perpendiculaire au plan (BDG) passant par L.
	\begin{enumerate}
		\item Justifier qu'une représentation paramétrique de la droite $\Delta$ est : 
\renewcommand\arraystretch{1.9}
\[
\left\{\begin{array}{l c l}x&=&\dfrac78 + t \\ y &=& \dfrac78+t \\ z&=&\dfrac34 -t \end{array}\right. \text{ où } t \in \R. \]
\renewcommand\arraystretch{1}
		\item Montrer que les droites $\Delta$ et (AE) sont sécantes au point K de coordonnées $\left(0~;~0~;~\dfrac{13}{8}\right)$.
		\item Que représente le point L pour le point K ? Justifier la réponse.
	\end{enumerate}
\item 
	\begin{enumerate}
		\item Calculer la distance KL.
		\item On admet que le triangle DBG est équilatéral. 
		
		Montrer que son aire est égale à $\dfrac{\sqrt{3}}{2}$.

		\item En déduire le volume du tétraèdre KDBG.

\begin{list}{\textbullet}{On rappelle que :}
\item le volume d'une pyramide est donné par la formule $\mathcal{V} = \dfrac13 \times \mathcal{B} \times h$ où $\mathcal{B}$ est l'aire d'une base et $h$ la longueur de la hauteur relative a cette base ;
\item un tétraèdre est une pyramide à base triangulaire.
\end{list}
	\end{enumerate}
	
\item On désigne par $a$ un réel appartenant à l'intervalle $]0~;~+\infty[$ et on note $K_a$ le point de coordonnées $(0~;~0~;~a)$.
	\begin{enumerate}
		\item Exprimer le volume $\mathcal{V}_a$ de la pyramide ABCD$K_a$ en fonction de $a$.
		\item On note $\Delta_a$ la droite de représentation paramétrique 

\[
\left\{\begin{array}{l c l}x&=&\phantom{-}t' \\ y &=&\phantom{-}t' \\ z&=&- t' + a \end{array}\right. \text{ où } t' \in \R. \]

On appelle $L_a$ le point d'intersection de la droite $\Delta_a$ avec le plan (BDG).

Montrer que les coordonnées du point $L_a$ sont $\left(\dfrac{a+1}{3}~;~\dfrac{a+1}{3}~;~\dfrac{2a-1}{3}\right)$.
		\item Déterminer, s'il existe, un réel strictement positif $a$ tel que le tétraèdre GDB$K_a$ et la pyramide ABCD$K_a$ sont de même volume.
	\end{enumerate}
\end{enumerate}

\end{document}
